\documentclass[UTF8,a4paper,12pt]{ctexart}

\usepackage{amsmath,amssymb}
\usepackage{booktabs}
\usepackage{longtable}
\usepackage{array}
\usepackage{geometry}
\usepackage{enumitem}
\usepackage{float}
\usepackage{multirow}
\usepackage{hyperref}
\usepackage{xcolor}
\usepackage[most]{tcolorbox}

\geometry{left=2.3cm,right=2.3cm,top=2.5cm,bottom=2.5cm}
\renewcommand{\arraystretch}{1.18}
\setlist[enumerate]{itemsep=2pt,topsep=2pt}
\setlist[itemize]{itemsep=2pt,topsep=2pt}

% 自定义修订说明框
\newtcolorbox{revisionbox}[1][]{
    colback=red!5!white,
    colframe=red!75!black,
    fonttitle=\bfseries,
    title=建模手v2.0核心修订与漏洞修复,
    #1
}

\title{\textbf{模型建立与求解方案 (v2.0 完整版)}}
\author{B题：可持续旅游管理}
\date{}

\begin{document}

\maketitle
\tableofcontents
\newpage

\section*{文档说明与版本更新}
\addcontentsline{toc}{section}{文档说明与版本更新}

本模块基于前文最终选择：小问1、2、3采用“方案2精简版”冲创新，小问4采用“方案1”保稳健，小问5采用“方案1增强版”保可解并兼顾空间均衡。数据口径统一采用赛题主口径：朱诺常住人口约3万人、2023年邮轮旅客160万人、旅游收入3.75亿美元、日峰值约2万人、港口物理容量7艘/日、政策上限5艘/日、邮轮季4—10月。

在初版模型交由编码手求解时，系统暴露出求解器为追求数学最优而“钻空子”的现象（模型漏洞）。本 v2.0 版由建模组进行了严密修正，使模型不仅能够收敛，且深度契合现实业务逻辑。

\begin{revisionbox}
\textbf{v2.0 核心修订项：}
\begin{enumerate}
    \item \textbf{收入底线“贴底走”修正}：v1.0中求解器为满足环保和限流，将收入压至 3.45 亿。v2.0强制将 \(R_{\min}\) 设为 3.6 亿美元，逼迫模型通过“提税（\(\tau\)）”而非“赶客”维持收入。
    \item \textbf{参数与数据口径一致性}：废弃主观设定的默认值。生态阈值 \(E_{\max}\) 严格取预处理数据上限 \(0.95\)；居民接受度下限 \(S_{\min}\) 依据调查放宽至 \(0.70\)。
    \item \textbf{封堵“生态作弊”漏洞}：原公式中环境成本可被无限抵扣，且游客量不受环境质量反噬。v2.0在需求函数引入环境吸引力因子 \((C_{env,0}/C_{env})^\gamma\)，并给环境成本设定物理下限，彻底封杀算法投机空间。
    \item \textbf{“非紧约束”学术化}：针对生态敏感性为 0 的现象，从运筹学 KKT 条件出发解释为“非活跃约束”，转化为政策建议。
\end{enumerate}
\end{revisionbox}

---

\section{全局符号、集合与数据口径}

\subsection{集合与下标}

\begin{longtable}{@{}p{2.2cm}p{4.0cm}p{8.5cm}@{}}
\caption{集合与下标定义}\\
\toprule
符号 & 含义 & 说明\\
\midrule
\endfirsthead
\toprule
符号 & 含义 & 说明\\
\midrule
\endhead
\(j\in\mathcal J\) & 管控措施编号 & \(j=1\)日上限，2酒类限制，3岸电，4交通管理，5营销分流\\
\(k\in\mathcal K\) & 支出项目编号 & 生态保护、基础设施、社区项目\\
\(s\in\mathcal S\) & 情景编号 & 低/中/高游客量、税率扰动、生态阈值等，来自预处理15情景\\
\(i\in\mathcal I\) & 区域编号 & 港口/市中心、门登霍尔冰川、观鲸、雨林等\\
\(r\in\mathcal R\) & 敏感性/措施组合情景 & 0-1措施组合枚举结果\\
\(t\in\mathcal T\) & 季节/月份 & 邮轮季4—10月\\
\bottomrule
\end{longtable}

\subsection{核心参数表 (v2.0 修正)}

\begin{longtable}{@{}p{1.8cm}p{2.4cm}p{3.2cm}p{1.5cm}p{5.5cm}@{}}
\caption{核心参数表}\\
\toprule
符号 & 类型/单位 & 数值或范围 & 属性 & 说明\\
\midrule
\endfirsthead
\toprule
符号 & 类型/单位 & 数值或范围 & 属性 & 说明\\
\midrule
\endhead
\(P\) & 参数/万人 & 3 & 已知 & 朱诺常住人口\\
\(x_0\) & 参数/万人 & 160 & 已知 & 2023邮轮旅客基准量\\
\(R_0\) & 参数/亿美元 & 3.75 & 已知 & 旅游收入基准\\
\(p_0\) & 参数/亿美元·万人\(^{-1}\) & \(3.75/160=0.0234375\) & 已知 & 人均旅游支出\\
\(x_{d0}\) & 参数/万人·日\(^{-1}\) & 2 & 已知 & 最繁忙日游客峰值\\
\(B_{\mathrm{phys}}\) & 参数/艘·日\(^{-1}\) & 7 & 已知 & 港口物理停靠容量\\
\(B_{\mathrm{pol}}\) & 参数/艘·日\(^{-1}\) & 5 & 已知 & 已实施政策上限\\
\(\tau_0\) & 参数/无量纲 & 0.11 & 已知 & 基准综合税率/价格杠杆\\
\(\tau_{\max}\) & 参数/无量纲 & 0.15 & 假设 & 税率上限\\
\(\varepsilon\) & 参数/无量纲 & \(-0.5\) & 假设 & 需求价格弹性区间\\
\(\alpha,\beta\) & 参数/无量纲 & \(0.2\) & 假设 & 容量、社会接受度对需求的正向弹性\\
\(\gamma\) & 参数/无量纲 & \textbf{0.1} & \textbf{新增} & \textbf{环境吸引力弹性系数}\\
\(\eta_I,\eta_E,\eta_C\) & 参数/无量纲 & \([0.5,0.9]\) & 假设 & 基础设施、生态、社区投资边际效率\\
\(\delta_{\mathrm{total}}\) & 参数/无量纲 & 基准0.8，上限1 & 已知/假设 & 新增税收可用于支出的总比例\\
\(E_{\max}\) & 参数/无量纲 & \textbf{0.95} & \textbf{修正} & \textbf{生态压力阈值（对齐真实数据）}\\
\(S_0\) & 参数/无量纲 & 1.0 & 假设 & 居民接受度基准指数\\
\(S_{\min}\) & 参数/无量纲 & \textbf{0.70} & \textbf{修正} & \textbf{居民接受度软约束下限}\\
\(R_{\min}\) & 参数/亿美元 & \textbf{3.6} & \textbf{修正} & \textbf{旅游收入下限，保障生计}\\
\(C_{env}^{\min}\) & 参数/无量纲 & \textbf{0.5} & \textbf{新增} & \textbf{基础生态压力防作弊下限}\\
\(B_0\) & 参数/亿美元 & \(\tau_0p_0x_0=0.4125\) & 推导 & 基准新增税收规模\\
\(c_k\) & 参数/亿美元 & 项目预算表 & 预处理 & 项目\(k\)成本\\
\(w_{ij}^{\mathrm{edge}}\) & 参数/无量纲 & \(0.6\tilde t_{ij}+0.4/\tilde{cap}_{ij}\) & 预处理 & 网络边权\\
\bottomrule
\end{longtable}

\subsection{全局变量表}

\begin{longtable}{@{}p{1.5cm}p{1.4cm}p{1.2cm}p{2.1cm}p{1.2cm}p{5.0cm}@{}}
\caption{全局变量表}\\
\toprule
符号 & 类型 & 单位 & 范围 & 属性 & 决策意义\\
\midrule
\endfirsthead
\toprule
符号 & 类型 & 单位 & 范围 & 属性 & 决策意义\\
\midrule
\endhead
\(x\) & 连续 & 万人/年 & \([0,I_{cap}]\) & 优化变量 & 年游客总量，核心规模控制变量\\
\(x_d\) & 连续 & 万人/日 & \([0,2]\) & 优化变量 & 日峰值游客量，拥堵与港口压力变量\\
\(\tau\) & 连续 & 无量纲 & \([\tau_{\min},\tau_{\max}]\) & 优化变量 & 综合税率/游客费价格杠杆\\
\(\delta_I,\delta_E,\delta_C\) & 连续 & 无量纲 & \([0,1]\)，和\(\le\delta_{\mathrm{total}}\) & 优化变量 & 基础设施、生态、社区支出比例\\
\(m_j\) & 0-1 & — & \(\{0,1\}\) & 优化变量 & \(m_j=1\)表示采用第\(j\)类管控措施\\
\(y_k\) & 0-1 & — & \(\{0,1\}\) & 优化变量 & \(y_k=1\)表示选择第\(k\)个支出项目\\
\(z_i\) & 0-1 & — & \(\{0,1\}\) & 优化变量 & \(z_i=1\)表示区域\(i\)热点限流\\
\(v_i\) & 0-1 & — & \(\{0,1\}\) & 优化变量 & \(v_i=1\)表示区域\(i\)冷点投资/营销\\
\(\lambda_{i,r}\) & 0-1 & — & \(\{0,1\}\) & 优化变量 & 区域\(i\)税率弹性分段线性化选择第\(r\)段\\
\(R\) & 中间 & 亿美元 & — & 目标变量 & 旅游总收入\\
\(B\) & 中间 & 亿美元 & — & 目标变量 & 新增税收收入\\
\(I_{cap}\) & 中间 & 万人/年 & — & 中间变量 & 基础设施容量\\
\(C_{env}\) & 中间 & 无量纲 & — & 目标变量 & 环境压力/生态成本指数\\
\(C_{soc}\) & 中间 & 无量纲 & — & 目标变量 & 社会成本/拥堵噪声压力指数\\
\(C_{inf}\) & 中间 & 无量纲 & — & 目标变量 & 基础设施压力指数\\
\(S\) & 中间 & 无量纲 & \([0,1]\) & 中间变量 & 居民接受度指数\\
\(D(\cdot)\) & 函数 & 万人/年 & — & 中间 & 需求响应函数\\
\(J\) & 目标 & 亿美元当量 & — & 目标变量 & 综合可持续效益\\
\bottomrule
\end{longtable}

\section{核心假设条件}

\begin{enumerate}[label=\textbf{A\arabic*.},leftmargin=2.2em]
\item \textbf{主口径一致性假设。} 模型主计算仅使用题给160万人、3.75亿美元、2万峰值、3万人口、7艘物理/5艘政策、4—10月邮轮季口径；参考PDF的167万、4.9亿美元、3850就业仅作为 \texttt{ref\_} 校准与讨论，不进入主模型。合理性：避免统计口径混算，是美赛数据规范的重要得分点。
\item \textbf{年聚合+日峰值硬约束假设。} 年游客量\(x\)与日峰值\(x_d\)分别建模，季节性通过4—10月容量与需求乘子体现，不逐日仿真。合理性：题目核心是年度可持续规模与最繁忙日拥挤，日粒度全仿真缺少数据且会显著增加不可解风险。
\item \textbf{需求弹性区间假设。} 需求对税率、容量、社会接受度采用对数弹性形式，弹性参数取区间\([-0.8,-0.2]\)及\([0.1,0.4]\)，并用分段线性/情景枚举校准。合理性：题给数据不足以精确估计弹性，区间情景比单点估计更稳健。
\item \textbf{支出反哺边际效率阶梯假设。} 新增税收投入基础设施、生态、社区后，容量、环境压力、居民接受度按边际效率\(\eta\in[0.5,0.9]\)三档阶梯提升；项目未选中\(y_k=0\)，未实施措施\(m_j=0\)。合理性：符合预处理“未实施=0”逻辑，避免“投一点钱无限扩容”的失真。
\item \textbf{区域功能分区与网络容量等级假设。} 将朱诺抽象为港口/市中心、冰川、观鲸、雨林等4—6个功能节点，网络边容量用等级参数和预处理边权\(w_{ij}^{\mathrm{edge}}=0.6\tilde t_{ij}+0.4/\tilde{cap}_{ij}\)刻画。合理性：无细分路网数据，功能分区可保留核心空间均衡逻辑。
\item \textbf{约束优先级假设。} 硬约束冲突时按“生态阈值 \(>\) 容量/5艘政策 \(>\) 预算平衡 \(>\) 居民接受度 \(>\) 收入下限”分层松弛，软约束用罚函数。合理性：题目强调冰川不可逆与过度旅游，生态和容量必须优先于收入。
\end{enumerate}

\section{小问1：朱诺可持续旅游基础模型 M1}

\subsection{变量定义与离散决策逻辑}

M1采用“多目标MINLP + 措施情景枚举 + NSGA-II/III + AHP-熵权折中”。连续变量控制规模与价格，0-1变量控制措施组合。

\begin{longtable}{@{}p{1.5cm}p{1.3cm}p{1.1cm}p{2.0cm}p{1.2cm}p{5.3cm}@{}}
\caption{M1变量定义}\\
\toprule
符号 & 类型 & 单位 & 范围 & 属性 & 决策意义\\
\midrule
\endfirsthead
\toprule
符号 & 类型 & 单位 & 范围 & 属性 & 决策意义\\
\midrule
\endhead
\(x\) & 连续 & 万人/年 & \([0,I_{cap}]\) & 优化 & 年游客量\\
\(x_d\) & 连续 & 万人/日 & \([0,2]\) & 优化 & 日峰值\\
\(\tau\) & 连续 & 无量纲 & \([0.09,0.15]\) & 优化 & 综合税率/游客费\\
\(\delta_I,\delta_E,\delta_C\) & 连续 & 无量纲 & \([0,1]\)，和\(\le\delta_{\mathrm{total}}\) & 优化 & 三类支出比例\\
\(m_1\) & 0-1 & — & \(\{0,1\}\) & 优化 & 日上限措施；\(m_1=1\)时压低\(x_d\)上限\\
\(m_2\) & 0-1 & — & \(\{0,1\}\) & 优化 & 酒类限制；\(m_2=1\)降低噪声社会成本\\
\(m_3\) & 0-1 & — & \(\{0,1\}\) & 优化 & 岸电；\(m_3=1\)降低环境压力\\
\(m_4\) & 0-1 & — & \(\{0,1\}\) & 优化 & 交通管理；\(m_4=1\)降低拥堵社会成本\\
\(m_5\) & 0-1 & — & \(\{0,1\}\) & 优化 & 营销分流；\(m_5=1\)提升需求/冷点分流\\
\(I_{cap}\) & 中间 & 万人/年 & — & 中间 & 基础设施容量\\
\(B\) & 中间 & 亿美元 & — & 中间 & 新增税收\\
\(S\) & 中间 & 无量纲 & \([0,1]\) & 中间 & 居民接受度\\
\bottomrule
\end{longtable}

\textbf{离散—公式强关联：} \(m_1\)进入日峰值上限式(M1-8)；\(m_2,m_4\)进入社会成本式(M1-5)；\(m_3\)进入环境压力式(M1-4')；\(m_5\)进入需求乘子式(M1-2')。因此每个0-1措施都对应可解释的约束或成本通道。

\subsection{公式推导 (v2.0 修正)}

\textbf{第1步：收入与税收。} 
\begin{equation}\tag{M1-1}
R=p_0x,\qquad B=\tau p_0x=\tau R.
\end{equation}
决策逻辑：式(M1-1)将游客量\(x\)与税率\(\tau\)关联到收入\(R\)和新增税收\(B\)。

\textbf{第2步：需求响应与漏洞封堵。} 为符合“冰川消失，游客终将不复存在”的题意，v2.0 在需求函数中引入环境吸引力因子：
\begin{equation}\tag{M1-2'}
x=x_0\left(\frac{\tau}{\tau_0}\right)^\varepsilon
\left(\frac{I_{cap}}{I_0}\right)^\alpha
\left(\frac{S}{S_0}\right)^\beta
\left(\frac{C_{env,0}}{C_{env} + 0.1}\right)^\gamma
\prod_{j=1}^{5}(1+\kappa_jm_j).
\end{equation}
决策逻辑：\(\gamma>0\)。若优化器试图通过压缩游客来降低环境成本，下一轮迭代中游客会因环境变好而反弹，倒逼模型限制过度的生态投资。

\textbf{第3步：容量反馈。}
\begin{equation}\tag{M1-3}
I_{cap}=I_0\left(1+\eta_I\delta_I\frac{B}{B_0}\right),\qquad B_0=\tau_0p_0x_0.
\end{equation}
决策逻辑：\(\delta_I\)为支出比例，\(\eta_I\)为边际效率；只有选择基础设施支出比例，容量才提升。

\textbf{第4步：环境成本抗作弊下限。}
\begin{equation}\tag{M1-4'}
C_{env} = \max \left\{ C_{env}^{\min}, \ e_1\frac{x}{x_0} + e_2\frac{x_d}{2}(1-\kappa_{\mathrm{power}}m_3) - \eta_E\delta_E\frac{B}{B_0} \right\}
\end{equation}
决策逻辑：设 \(C_{env}^{\min} = 0.5\)。只要有人类活动就存在基础环境压力，生态投资只能减轻污染，不能将其完全抹零。

\textbf{第5步：社会成本。}
\begin{equation}\tag{M1-5}
C_{soc}=s_h\left(\frac{x}{P}\right)^2
+s_n\frac{x_d}{2}(1-\kappa_{\mathrm{traffic}}m_4)
+s_{\mathrm{noise}}(1-\kappa_{\mathrm{alc}}m_2).
\end{equation}
决策逻辑：游客/居民比平方项刻画拥挤非线性；\(m_4\)交通管理降低拥堵；\(m_2\)酒类限制降低噪声。

\textbf{第6步：基础设施压力。}
\begin{equation}\tag{M1-6}
C_{inf}=i_p\left(\frac{x}{I_{cap}}\right)^2+i_d\frac{x_d}{2}.
\end{equation}
决策逻辑：容量利用率平方项体现接近容量时的非线性压力。

\textbf{第7步：居民接受度。}
\begin{equation}\tag{M1-7}
S=S_0-s_x\frac{x}{x_0}
+s_d\frac{x_d}{2}(1-\kappa_{\mathrm{traffic}}m_4)
+\eta_C\delta_C\frac{B}{B_0}
+\kappa_{\mathrm{mark}}m_5.
\end{equation}
决策逻辑：交通管理和社区支出提升接受度，营销分流\(m_5\)改善冷点体验。

\textbf{第8步：硬约束。}
\begin{equation}\tag{M1-8}
0\le x,\qquad 0\le x_d\le 2-\Delta_{\mathrm{cap}}m_1,\qquad \text{船数/日}\le 5\le 7.
\end{equation}
\begin{equation}\tag{M1-9}
x\le I_{cap},\qquad C_{env}\le 0.95,\qquad 0.09\le\tau\le0.15.
\end{equation}
\begin{equation}\tag{M1-10}
\delta_I+\delta_E+\delta_C\le\delta_{\mathrm{total}}\le1,\qquad m_j\in\{0,1\}.
\end{equation}
决策逻辑：v2.0将生态阈值严格设为 0.95；5艘/日政策上限与7艘物理容量分开。

\textbf{第9步：软约束与罚项。}
\begin{equation}\tag{M1-11}
P_{\mathrm{soft}}=\lambda_S\max(0,0.70-S)+\lambda_R\max(0,3.6-R)+\lambda_H\max(0,C_{soc}-C_{soc}^{\max}).
\end{equation}
决策逻辑：v2.0强制拉升收入底线至 3.6 亿，逼迫模型用提税来弥补限流损失，不能单纯压收入。

\textbf{第10步：多目标综合。}
\begin{equation}\tag{M1-12}
\begin{aligned}
\max&\ f_1=R, & \min&\ f_2=C_{env},\\
\min&\ f_3=C_{soc}, & \min&\ f_4=C_{inf},
\end{aligned}
\end{equation}
\begin{equation}\tag{M1-13}
J=w_1\frac{f_1}{R_0}-w_2f_2-w_3f_3-w_4f_4-P_{\mathrm{soft}}.
\end{equation}

\textbf{第11步：AHP-熵权折中。}
\begin{equation}\tag{M1-14}
p_{ij}=\frac{f_{ij}}{\sum_i f_{ij}},\qquad
e_j=-\frac{1}{\ln n}\sum_i p_{ij}\ln p_{ij},\qquad
w_j^{\mathrm{ent}}=\frac{1-e_j}{\sum_j(1-e_j)}.
\end{equation}
AHP权重\(w_j^{\mathrm{AHP}}\)由判断矩阵最大特征向量归一化；组合权重为：
\begin{equation}\tag{M1-15}
w_j=\theta w_j^{\mathrm{AHP}}+(1-\theta)w_j^{\mathrm{ent}},\qquad 0\le\theta\le1.
\end{equation}

\subsection{M1建模流程图文字描述}

决策场景抽象（3万人/160万游客/3.75亿美元/2万峰值/7艘物理+5艘政策/4—10月）\(\rightarrow\) 定义连续变量\(x,x_d,\tau,\delta\)与0-1措施\(m_j\) \(\rightarrow\) 建立需求响应式(M1-2')与容量反馈式(M1-3) \(\rightarrow\) 建立环境/社会/基础设施成本式(M1-4')—(M1-6) \(\rightarrow\) 建立居民接受度式(M1-7) \(\rightarrow\) 列硬约束(M1-8)—(M1-10)与软约束(M1-11) \(\rightarrow\) \textbf{关键节点：检查\(5\le7\)、\(\tau_{\min}\le\tau_{\max}\)、\(\delta\)和不超过\(\delta_{\mathrm{total}}\)、\(x\le I_{cap}\)是否兼容；若冲突按生态\(>\)容量\(>\)预算\(>\)居民\(>\)收入松弛} \(\rightarrow\) 对\(m_j\)做0-1情景枚举/混合编码 \(\rightarrow\) 采用不动点迭代法处理\(x\)与\(I_{cap}\)的循环依赖 \(\rightarrow\) NSGA-II/III迭代 \(\rightarrow\) 生成Pareto前沿 \(\rightarrow\) AHP-熵权选折中 \(\rightarrow\) 输出变量表、约束表、基准方案、措施优先级、Pareto图。

---

\section{小问2：新增税收支出与反哺模型 M2}

\subsection{变量定义}

\begin{longtable}{@{}p{1.5cm}p{1.3cm}p{1.1cm}p{2.2cm}p{1.2cm}p{5.1cm}@{}}
\caption{M2变量定义}\\
\toprule
符号 & 类型 & 单位 & 范围 & 属性 & 决策意义\\
\midrule
\endfirsthead
\toprule
符号 & 类型 & 单位 & 范围 & 属性 & 决策意义\\
\midrule
\endhead
\(\delta_I,\delta_E,\delta_C\) & 连续 & 无量纲 & \([0,1]\)，和\(\le\delta_{\mathrm{total}}\) & 优化 & 三类支出比例\\
\(y_k\) & 0-1 & — & \(\{0,1\}\) & 优化 & \(y_k=1\)选择项目\(k\)，进入预算与反哺效果\\
\(s\in\mathcal S\) & 情景 & — & 低/中/高 & 参数 & 游客量、税率、收入不确定情景\\
\(B_s\) & 中间 & 亿美元 & — & 中间 & 情景\(s\)下新增税收\\
\(\Delta I_s,\Delta E_s,\Delta S_s\) & 中间 & 无量纲 & — & 中间 & 基础设施、生态、社会反哺效果\\
\(U_s\) & 中间 & 亿美元当量 & — & 目标 & 情景\(s\)综合效用\\
\(z\) & 连续 & 亿美元当量 & — & 目标 & 鲁棒最坏情景效用\\
\bottomrule
\end{longtable}

\subsection{公式推导}

\textbf{第1步：情景税收。}
\begin{equation}\tag{M2-1}
B_s=\tau p_0x_s,\qquad B_{\min}=\min_{s\in\mathcal S}B_s.
\end{equation}

\textbf{第2步：鲁棒预算约束。}
\begin{equation}\tag{M2-2}
\sum_{k\in\mathcal K}c_ky_k\le\delta_{\mathrm{total}}B_{\min},\qquad \delta_I+\delta_E+\delta_C\le\delta_{\mathrm{total}}.
\end{equation}

\textbf{第3步：项目反哺效果。}
\begin{equation}\tag{M2-3}
\Delta I_s=\sum_k\eta_{Ik}\frac{c_ky_k}{B_0},\quad
\Delta E_s=\sum_k\eta_{Ek}\frac{c_ky_k}{B_0},\quad
\Delta S_s=\sum_k\eta_{Ck}\frac{c_ky_k}{B_0}.
\end{equation}

\textbf{第4步：反馈闭环。}
\begin{equation}\tag{M2-4}
I_s=I_0(1+\Delta I_s),\quad
C_{env,s}=\max(0.5, C_{env,0}(1-\rho_E\Delta E_s)),\quad
S_s=S_0+\Delta S_s,
\end{equation}
\begin{equation}\tag{M2-5'}
x_s^{\mathrm{new}}=x_0\left(\frac{\tau}{\tau_0}\right)^\varepsilon
\left(\frac{I_s}{I_0}\right)^\alpha
\left(\frac{S_s}{S_0}\right)^\beta
\left(\frac{C_{env,0}}{C_{env,s} + 0.1}\right)^\gamma
\prod_j(1+\kappa_jm_j).
\end{equation}
决策逻辑：v2.0在此处同样应用了环境反噬与防作弊下限，确保反馈闭环逻辑一致。

\textbf{第5步：情景效用。}
\begin{equation}\tag{M2-6}
U_s=\lambda_I\Delta I_s+\lambda_E\Delta E_s+\lambda_C\Delta S_s
-\lambda_L\max(0,x_0-x_s^{\mathrm{new}})
-\lambda_B\max(0,C_{env,s}-0.95).
\end{equation}

\textbf{第6步：鲁棒目标。}
\begin{equation}\tag{M2-7}
\max z \quad \text{s.t.}\quad z\le U_s,\ \forall s\in\mathcal S,
\end{equation}

\textbf{第7步：项目互斥与逻辑约束。}
\begin{equation}\tag{M2-9}
y_a+y_b\le1,\quad (a,b)\in\mathcal M_{\mathrm{excl}},\qquad y_k\in\{0,1\}.
\end{equation}

\subsection{M2建模流程图文字描述}

识别生态/基础设施/社区项目 \(\rightarrow\) 设定预算上限\(B_{\min}\)与互斥规则 \(\rightarrow\) 定义0-1项目\(y_k\)和连续比例\(\delta\) \(\rightarrow\) 建立预算式(M2-2) \(\rightarrow\) 建立反哺效果式(M2-3) \(\rightarrow\) SD三回路反馈式(M2-4)—(M2-5') \(\rightarrow\) 蒙特卡洛生成低/中/高情景 \(\rightarrow\) \textbf{关键节点：检查最坏情景预算是否可行；若不可行，削减低边际效率项目或增加保守储备} \(\rightarrow\) 计算情景效用(M2-6) \(\rightarrow\) 鲁棒目标(M2-7) \(\rightarrow\) 输出支出分配表、项目清单、闭环反馈图、情景对比与置信区间。

---

\section{小问3：敏感性分析模型 M3}

\subsection{变量定义}

\begin{longtable}{@{}p{1.5cm}p{1.4cm}p{1.2cm}p{2.2cm}p{1.2cm}p{5.0cm}@{}}
\caption{M3变量定义}\\
\toprule
符号 & 类型 & 单位 & 范围 & 属性 & 决策意义\\
\midrule
\endfirsthead
\toprule
符号 & 类型 & 单位 & 范围 & 属性 & 决策意义\\
\midrule
\endhead
\(\theta_p\) & 参数 & — & 扰动区间 & 优化/情景 & 待敏感性参数：\(\varepsilon,\tau_{\max},E_{\max},S_{\min},\eta_I,\eta_E,\eta_C,w_j\)\\
\(m_r\) & 0-1向量 & — & 枚举 & 离散情景 & 第\(r\)个措施组合\\
\(S_p\) & 中间 & 无量纲 & — & 目标 & 单因素弹性系数\\
\(S_i,S_{Ti}\) & 中间 & 无量纲 & \([0,1]\) & 目标 & Sobol一阶/总效应指数\\
\(B_r\) & 中间 & 分 & — & 目标 & 措施组合Borda稳健得分\\
\(\theta^{\mathrm{break}}\) & 中间 & 参数单位 & — & 目标 & 硬约束破坏阈值\\
\bottomrule
\end{longtable}

\subsection{公式推导与学术发现}

\textbf{第1步：单因素弹性。}
\begin{equation}\tag{M3-1}
S_p=\frac{\Delta J/J}{\Delta p/p}.
\end{equation}

\textbf{第2步：Sobol全局敏感性。}
\begin{equation}\tag{M3-2}
S_i=\frac{\operatorname{Var}_{\theta_i}\left(\mathbb E[Y|\theta_i]\right)}{\operatorname{Var}(Y)},\qquad S_{Ti}=1-\frac{\operatorname{Var}_{\theta_{-i}}\left(\mathbb E[Y|\theta_{-i}]\right)}{\operatorname{Var}(Y)}.
\end{equation}

\begin{revisionbox}[title=v2.0 核心学术发现：生态非紧约束]
在 M1 v2.0 的基准求解中，最优解的生态压力 \(C_{env} \approx 0.5\)（刚好压在防作弊下限），远低于红线 \(E_{\max}=0.95\)。因此扰动 \(E_{\max}\) 时，目标函数的弹性系数 \(S_p\) 均为 0。

\textbf{数学解释}：由于约束 \(C_{env} \le E_{\max}\) 存在极大松弛量（Slack），在 KKT 条件下该约束为\textbf{非活跃约束（Inactive Constraint）}，其影子价格为 0。
\textbf{业务结论}：这说明当前的管控组合（如岸电 \(m_3\) 和适度限流 \(m_1\)）已经极为有效，旅游业的局部环境压力已被成功遏制。朱诺下一步的治理重心应从“单纯限流保冰川”，转向“通过税收反哺缓解居民社会压力（提升 \(S\)）”。
\end{revisionbox}

\textbf{第3步：硬约束破坏阈值。}
\begin{equation}\tag{M3-3}
\theta_p^{\mathrm{break}}=\inf\{\theta_p:\ \text{M1/M2硬约束不可行}\}.
\end{equation}

\textbf{第4步：离散措施组合秩和稳健性。}
\begin{equation}\tag{M3-4}
R_{rq}=\operatorname{rank}\left(J_r(\theta_q)\right),\qquad B_r=\sum_q\left(N-R_{rq}+1\right).
\end{equation}

\textbf{第5步：稳健措施选择。}
\begin{equation}\tag{M3-5}
\max_r B_r\quad \text{s.t.}\quad \mathrm{Feasible}_r(\theta_q)=1,\ \forall q.
\end{equation}

\subsection{M3建模流程图文字描述}

选定参数与离散措施情景 \(\rightarrow\) 单因素扰动/情景枚举 \(\rightarrow\) 重复求解M1/M2 \(\rightarrow\) \textbf{关键节点：检查硬约束是否可行；不可行则记录破坏阈值并触发优先级松弛} \(\rightarrow\) 计算弹性系数\(S_p\) \(\rightarrow\) 筛3—5个关键参数 \(\rightarrow\) 蒙特卡洛抽样 \(\rightarrow\) Sobol指数计算 \(\rightarrow\) \textbf{关键节点：若交互效应显著，增加参数组合情景；否则输出单参数阈值} \(\rightarrow\) 措施组合Borda排序 \(\rightarrow\) 输出关键因素排序表、弹性系数表、Sobol指数图、阈值图、稳健措施组合清单。

---

\section{小问4：迁移到另一过度旅游目的地 M4}

\subsection{变量定义}

选择\textbf{威尼斯}作为迁移对象，与朱诺形成“自然生态型 vs 城市文化遗产型”对比。

\begin{longtable}{@{}p{1.5cm}p{1.3cm}p{1.2cm}p{2.2cm}p{1.2cm}p{5.1cm}@{}}
\caption{M4变量定义}\\
\toprule
符号 & 类型 & 单位 & 范围 & 属性 & 决策意义\\
\midrule
\endfirsthead
\toprule
符号 & 类型 & 单位 & 范围 & 属性 & 决策意义\\
\midrule
\endhead
\(d\) & 离散 & — & \(\{J,V\}\) & 优化/情景 & 目的地选择，\(J\)朱诺，\(V\)威尼斯\\
\(F_k^d\) & 参数 & 无量纲 & \([0,1]\) & 已知/假设 & 目的地特征：自然依赖、城市密度、住房压力、交通拥堵、生态敏感、季节集中、居民反对、管理能力\\
\(w_k^d\) & 连续 & 无量纲 & \([0,1]\)，和1 & 优化/中间 & 目的地权重\\
\(m_j^d\) & 0-1 & — & \(\{0,1\}\) & 优化 & 目的地\(d\)是否采用措施\(j\)\\
\(C_j^d\) & 中间 & 无量纲 & \([0,1]\) & 目标 & TOPSIS措施贴近度/优先级\\
\(\theta^d\) & 参数 & 原单位 & 重校准 & 中间 & 迁移后阈值/参数\\
\bottomrule
\end{longtable}

\subsection{公式推导}

\textbf{第1步：特征归一化。}
\begin{equation}\tag{M4-1}
\tilde F_k^d=\frac{F_k^d-\min_dF_k^d}{\max_dF_k^d-\min_dF_k^d}.
\end{equation}

\textbf{第2步：AHP权重。}
\begin{equation}\tag{M4-2}
A w^{\mathrm{AHP}}=\lambda_{\max}w^{\mathrm{AHP}},\qquad \sum_kw_k^{\mathrm{AHP}}=1.
\end{equation}

\textbf{第3步：熵权。}
\begin{equation}\tag{M4-3}
p_{kd}=\frac{\tilde F_k^d}{\sum_d\tilde F_k^d},\qquad e_k=-\frac{1}{\ln n}\sum_dp_{kd}\ln p_{kd},\qquad w_k^{\mathrm{ent}}=\frac{1-e_k}{\sum_k(1-e_k)}.
\end{equation}

\textbf{第4步：组合权重。}
\begin{equation}\tag{M4-4}
w_k^d=\theta w_k^{\mathrm{AHP}}+(1-\theta)w_k^{\mathrm{ent}}.
\end{equation}

\textbf{第5步：TOPSIS措施排序。}
\begin{equation}\tag{M4-5}
V_{jk}=w_k^d\tilde R_{jk},\qquad D_j^+=\sqrt{\sum_k(V_{jk}-V_k^+)^2},\qquad D_j^-=\sqrt{\sum_k(V_{jk}-V_k^-)^2},
\end{equation}
\begin{equation}\tag{M4-6}
C_j^d=\frac{D_j^-}{D_j^++D_j^-}.
\end{equation}

\textbf{第6步：阈值重校准。}
\begin{equation}\tag{M4-7}
\theta^d=\operatorname{clip}\left(\theta^J\left(1+\sum_k\gamma_{\theta k}\tilde F_k^d\right),\underline\theta,\bar\theta\right).
\end{equation}

\textbf{第7步：迁移优化模型。}
\begin{equation}\tag{M4-8}
\max J^d=w_1^d\frac{R^d}{R_0^d}-w_2^dC_{env}^d-w_3^dC_{soc}^d-w_4^dC_{inf}^d-P_{\mathrm{soft}}^d.
\end{equation}

\subsection{M4建模流程图文字描述}

选择迁移目的地威尼斯 \(\rightarrow\) 抽象特征向量并与朱诺对比 \(\rightarrow\) 归一化特征 \(\rightarrow\) AHP/熵权得组合权重 \(\rightarrow\) TOPSIS排序措施优先级 \(\rightarrow\) \textbf{关键节点：用式(M4-7)重校准阈值；若新目的地硬约束不可行，调整阈值或松弛软约束} \(\rightarrow\) 调用M1框架求解\(J^d\) \(\rightarrow\) 输出朱诺—威尼斯对比表、参数重校准表、措施优先级表、迁移优化结果。

---

\section{小问5：热门—冷门空间均衡模型 M5}

\subsection{变量定义}

\begin{longtable}{@{}p{1.5cm}p{1.3cm}p{1.2cm}p{2.2cm}p{1.2cm}p{5.1cm}@{}}
\caption{M5变量定义}\\
\toprule
符号 & 类型 & 单位 & 范围 & 属性 & 决策意义\\
\midrule
\endfirsthead
\toprule
符号 & 类型 & 单位 & 范围 & 属性 & 决策意义\\
\midrule
\endhead
\(x_i\) & 连续 & 万人/年 & \([0,cap_i]\) & 优化 & 区域\(i\)游客量\\
\(\tau_i\) & 连续 & 无量纲 & \([\tau_{\min},\tau_{\max}]\) & 优化 & 区域差异化税率\\
\(z_i\) & 0-1 & — & \(\{0,1\}\) & 优化 & \(z_i=1\)热点限流，0不限流\\
\(v_i\) & 0-1 & — & \(\{0,1\}\) & 优化 & \(v_i=1\)冷点投资/营销，0不投\\
\(\lambda_{i,r}\) & 0-1 & — & \(\{0,1\}\) & 优化 & 税率弹性分段线性化选择第\(r\)段\\
\(H_i\) & 中间 & 无量纲 & — & 目标 & 热点压力\(x_i/cap_i\)\\
\(U_i\) & 中间 & 无量纲 & — & 目标 & 冷点利用率\(x_i/cap_i\)\\
\(f_{ij}\) & 连续 & 万人/年 & \(\ge0\) & 中间 & 网络边\((i,j)\)游客流\\
\(cap_i\) & 参数 & 万人/年 & 等级 & 已知/假设 & 区域容量\\
\(\Delta cap_i\) & 参数 & 万人/年 & 等级 & 假设 & 冷点投资扩容\\
\(H_{\max}\) & 参数 & 无量纲 & 0.8—0.9 & 假设 & 热点压力阈值\\
\(\rho\) & 参数 & 无量纲 & 0.4—0.6 & 假设 & 冷点公平利用率下限\\
\bottomrule
\end{longtable}

\subsection{公式推导}

\textbf{第1步：空间收益目标。}
\begin{equation}\tag{M5-1}
\max \sum_i p_0(1+\tau_i)x_i -w_H\sum_i H_i +w_B\sum_i U_i -\sum_i c_i v_i.
\end{equation}

\textbf{第2步：游客守恒。}
\begin{equation}\tag{M5-2}
\sum_i x_i=x,\qquad x_i\ge0.
\end{equation}

\textbf{第3步：离散容量约束。}
\begin{equation}\tag{M5-3}
x_i\le cap_i(1-z_i)+\Delta cap_i v_i.
\end{equation}

\textbf{第4步：热点压力与冷点公平。}
\begin{equation}\tag{M5-4}
H_i=\frac{x_i}{cap_i}\le H_{\max},\qquad i\in\mathcal I_{\mathrm{hot}},
\end{equation}
\begin{equation}\tag{M5-5}
U_i=\frac{x_i}{cap_i}\ge \rho v_i,\qquad i\in\mathcal I_{\mathrm{cold}}.
\end{equation}

\textbf{第5步：税率弹性分段线性化。}
\begin{equation}\tag{M5-6}
x_i\le a_{i,r}-b_{i,r}\tau_i+M(1-\lambda_{i,r}),\qquad \sum_r\lambda_{i,r}=1,\qquad \lambda_{i,r}\in\{0,1\}.
\end{equation}

\textbf{第6步：简化网络流校验。}
\begin{equation}\tag{M5-7}
\sum_{i\in S}x_i\le C_{\delta(S)},\qquad \forall S\in\mathcal C,
\end{equation}

\textbf{第7步：投资预算与税率范围。}
\begin{equation}\tag{M5-8}
\sum_i c_i v_i\le B,\qquad \tau_{\min}\le\tau_i\le\tau_{\max}.
\end{equation}

\subsection{M5建模流程图文字描述}

区域划分与容量分级（港口/市中心热点、冰川热点、观鲸/雨林冷点）\(\rightarrow\) 定义\(x_i,\tau_i,z_i,v_i,\lambda_{i,r}\) \(\rightarrow\) 建立收益目标(M5-1) \(\rightarrow\) 游客守恒(M5-2) \(\rightarrow\) 离散容量式(M5-3) \(\rightarrow\) 热点压力/冷点公平(M5-4)(M5-5) \(\rightarrow\) 税率分段线性化(M5-6) \(\rightarrow\) 网络流割约束(M5-7) \(\rightarrow\) \textbf{关键节点：检查热点压力是否超阈值；若超，启用\(z_i=1\)或提高热点税率；检查冷点利用率是否低于公平下限；若低，增加\(v_i=1\)冷点投资/营销} \(\rightarrow\) MILP求解 \(\rightarrow\) 输出区域分流表、差异化税率表、投资表、热点压力对比图、均衡前后指标。

---

\section{建模流程图总链与论文写作建议}

总逻辑链可写为：
\[
\begin{aligned}
&\text{数据口径统一与预处理（160万/3.75亿/2万/3万/7艘物理+5艘政策/4—10月）}\\
&\rightarrow \text{M1基础多目标MINLP }(x,x_d,\tau,\delta,m_j)\\
&\rightarrow \text{M2税收反哺鲁棒MILP+SD三回路 }(y_k,\delta)\\
&\rightarrow \text{M3单因素+Sobol敏感性 }(\theta_p,m_r)\\
&\rightarrow \text{M4特征校准+AHP-TOPSIS+MINLP重校准 }(d,m_j^d,\theta^d)\\
&\rightarrow \text{M5容量分配MILP+公平约束+网络流校验 }(x_i,\tau_i,z_i,v_i,\lambda_{i,r})\\
&\rightarrow \text{备忘录汇总。}
\end{aligned}
\]

论文中建议每小问按“变量表\(\rightarrow\)假设\(\rightarrow\)公式推导\(\rightarrow\)离散逻辑说明\(\rightarrow\)约束兼容性检查\(\rightarrow\)流程图\(\rightarrow\)输出物”组织；所有公式编号采用(M1-1)、(M2-1)等；每个0-1变量至少在一处约束或目标中显式出现，避免“离散变量与公式脱节”的美赛常见扣分点。

\end{document}